\documentclass[11pt,a4paper]{article}
\usepackage[top=0.75in,bottom=0.75in,left=0.75in,right=0.75in]{geometry}
\usepackage[T1]{fontenc}
\usepackage{times,microtype,enumitem,hyperref,titlesec}
\pagestyle{empty}
\setlength{\parindent}{0pt}
\setlength{\parskip}{0pt}
\linespread{1.10}
\hypersetup{hidelinks,breaklinks=true}
\titleformat{\section}{\bfseries\normalsize}{}{0pt}{}[\titlerule]
\titlespacing*{\section}{0pt}{10pt}{4pt}
\newcommand{\entry}[3]{\noindent\textbf{#1}\hfill\textbf{#3}\\[-1pt]\textit{#2}\par}
\newenvironment{cvitems}{\begin{itemize}[leftmargin=1.1em,label=\textbullet,itemsep=2pt,topsep=3pt,parsep=0pt,partopsep=0pt]}{\end{itemize}}
\begin{document}
{\LARGE\bfseries Aarush Ram Anandh}\par\vspace{4pt}
\href{mailto:aarush.anandh_ug2024@ashoka.edu.in}{aarush.anandh\_ug2024@ashoka.edu.in}\quad|\quad +91 9659786783\par\vspace{2pt}
\href{https://github.com/MushuRam01}{github.com/MushuRam01}\quad|\quad
\href{https://www.linkedin.com/in/aarush-ram-anandh/}{linkedin.com/in/aarush-ram-anandh}\par
\vspace{20pt}
11 November 2026\par\vspace{12pt}
Hiring Team\\
Sony Research India\par\vspace{12pt}
\textbf{Application: LLM Research Internship}\par\vspace{14pt}
Dear Hiring Team,\par\vspace{10pt}
I am applying for Sony Research India\textquotesingle s LLM research internship. I am completing a B.Sc. (Hons.) in Computer Science and Mathematics at Ashoka University, expected in 2028. Your emphasis on understanding model behaviour and developing methods through hypotheses and experiments is particularly appealing to me. My hands-on work is in diffusion models, reinforcement learning, audio classification and Transformer-based financial modelling.\par\vspace{10pt}
At the Constrained Image-Synthesis Lab, I investigated inference-time colour preservation in latent diffusion, achieving $\Delta E_{00}<1$ and approximately four times better colour accuracy than evaluated open-source baselines. The work is reported in my public preprint, \textit{Inference-Time Loss-Guided Colour Preservation in Diffusion Sampling} (arXiv:2601.17259). I am also developing a controllable optimal-transport framework, CGSW, with a paper submitted to JMLR. These projects have involved studying why a method behaves as it does, translating mathematical ideas into code and choosing experiments that test the intended claim.\par\vspace{10pt}
In an ongoing two-person RL project, we co-developed masked-action PPO and DQN agents and evaluated PPO across two team variants and five seeds each. PPO reached 99.7--100\% wins against the fixed heuristic over 1,020 final battles per run; confidence intervals and permutation tests showed no significant learning-speed difference between variants. My MEIKURAL project classifies seven Hindustani raags with STFT/CNN representations trained on 3,471 audio windows; at the Koita Centre for Digital Health, I contributed to a wearable-data pipeline achieving approximately 70\% sensitivity for asthma exacerbation detection in study evaluation. I also built experimental frozen time-series foundation-model routes for surprise features and embeddings, implemented but not yet run in the reported cohort evaluation.\par\vspace{10pt}
At Avartan Labs, I raised rare-class correct detection from 40\% to 80\% in an RT-DETR pipeline covering 18 material classes. This industry internship gave me experience analysing detector failures and investigating targeted model improvements.\par\vspace{10pt}
I would welcome the opportunity to apply this experimental approach to multilingual model behaviour, reasoning and entertainment-domain adaptation. I can contribute Python/PyTorch prototyping and careful analysis while developing the language-model-specific expertise the work requires. Thank you for considering my application.\par\vspace{14pt}
Yours sincerely,\\[5pt]
Aarush Ram Anandh
\end{document}
